% Source: https://artofproblemsolving.com/community/c3h3113802_clue_combinatorics_conundrums
% Stable identity: aops-topic-3113802-q03
A matching pair exists in the dealt deck precisely when its original card was not hidden. These three independent probabilities are $\dfrac{5}{6},\dfrac{5}{6},\dfrac{8}{9}$. For $ j $ specified existing pairs, let $ T_j $ be the probability that each pair lies entirely in one of the two specified hands. Writing $(a)_k=a(a-1)\cdots(a-k+1)$, uniformly placing their $ 2j $ distinct physical cards gives \[T_j=\dfrac{\sum_{l=0}^j\binom jl(7)_{2l}(7)_{2j-2l}}{(21)_{2j}}.\]The values are $ T_1=\dfrac{1}{5},T_2=\dfrac{31}{855},T_3=\dfrac{11}{1938}$. The elementary symmetric sums of the three pair-existence probabilities are $\dfrac{23}{9},\dfrac{235}{108},\dfrac{50}{81}$. Inclusion-exclusion therefore gives \[\dfrac{23}{9}\dfrac{1}{5}-\dfrac{235}{108}\dfrac{31}{855}+\dfrac{50}{81}\dfrac{11}{1938}=\boxed{\dfrac{683987}{1569780}}.\]Only the three duplicated identities can produce a matching pair, so these events exhaust the possibilities.
